% theme: quadratic_equations_function_applications
\MajorQuestionLesson{関数の利用問題}
\SetRowSpace{6mm}

\TypeQuestionSingle{}{kansuu1}
\FigProblemBlock{45mm}{fig/lesson_10_kansuu_01}{\TextTall{点 ${\rm P}$ は関数 $y=\dfrac{1}{2}x+4$ のグラフ上にあり、点 ${\rm Q}$ は点 ${\rm P}$ から $x$ 軸に下ろした垂線の足である。点 ${\rm P}$ の $x$ 座標を $t$ とし、$t>0$ とする。}}
{
\QQ{$\Tri{OPQ}$ の面積を $t$ を使って表しなさい。}{}
\QQ{$\Tri{OPQ}$ の面積が21になるとき、$t$ の値を求めなさい。}{}
\QQ{$\Tri{OPQ}$ の面積が24になるとき、$t$ の値を求めなさい。}{}
}

\TypeQuestionSingle{}{kansuu2}
\FigProblemBlock{45mm}{fig/lesson_10_kansuu_02}{\TextTall{関数 $y=-x+12$ のグラフ上に点 ${\rm P}$ がある。点 ${\rm P}$ から $x$ 軸、$y$ 軸にそれぞれ垂線を下ろし、その足を ${\rm Q,R}$ とする。点 ${\rm P}$ の $x$ 座標を $t$ とし、$0<t<12$ とする。}}
{
\QQ{点 ${\rm P}$ の $y$ 座標を $t$ を使って表しなさい。}{}
\QQ{四角形 ${\rm OQPR}$ の面積を $t$ を使って表しなさい。}{}
\QQ{四角形 ${\rm OQPR}$ の面積が35になるとき、$t$ の値をすべて求めなさい。}{}
}

\LessonPageBreak

\TypeQuestionSingle{}{kansuu3}
\FigProblemBlock{45mm}{fig/lesson_10_kansuu_03}{\TextTall{2直線 $y=2x+1$ と $y=-x+13$ の交点を ${\rm R}$ とする。点 ${\rm P}$ は直線 $y=2x+1$ 上、点 ${\rm Q}$ は直線 $y=-x+13$ 上にあり、$\Lseg{PQ}$ は $y$ 軸に平行である。また、点 ${\rm P,Q}$ は点 ${\rm R}$ より右側にある。}}
{
\QQ{$\Tri{RPQ}$ の面積が6になるとき、点 ${\rm P}$ の座標を求めなさい。}{}
\QQ{$\Tri{RPQ}$ と $\Tri{OPQ}$ の面積の比が $1:2$ になるとき、点 ${\rm P}$ の座標を求めなさい。}{}
\QQ{$\Tri{RPQ}$ の面積が3になるとき、点 ${\rm P}$ の座標を求めなさい。}{}
}

\TypeQuestionSingle{}{kansuu4}
\FigProblemBlock{45mm}{fig/lesson_10_kansuu_04}{\TextTall{関数 $y=\dfrac{1}{3}x+2$ のグラフと $y$ 軸との交点を ${\rm R}$ とする。また、点 ${\rm P}$ はこのグラフ上にあり、点 ${\rm Q}$ は点 ${\rm P}$ から $x$ 軸に下ろした垂線の足である。点 ${\rm P}$ の $x$ 座標を $t$ とし、$t>0$ とする。4点 ${\rm O,Q,P,R}$ をこの順に結んだ台形を考える。}}
{
\QQ{台形 ${\rm OQPR}$ の面積が18になるとき、点 ${\rm P}$ の座標を求めなさい。}{}
\QQ{台形 ${\rm OQPR}$ の面積が $\Tri{OPQ}$ の面積の $\dfrac{4}{3}$ 倍になるとき、点 ${\rm P}$ の座標を求めなさい。}{}
\QQ{台形 ${\rm OQPR}$ の面積が15になるとき、点 ${\rm P}$ の座標を求めなさい。}{}
}
